% This is samplepaper.tex, a sample chapter demonstrating the
% LLNCS macro package for Springer Computer Science proceedings;
% Version 2.20 of 2017/10/04
%
\documentclass[runningheads]{llncs}
%
\usepackage{graphicx}
\usepackage{amssymb}
\usepackage{amsmath}
\usepackage{hyperref}
\usepackage{booktabs}
\usepackage{cleveref}
\usepackage{subcaption}
\hypersetup{pdfpagemode={UseOutlines},
bookmarksopen=true,
bookmarksopenlevel=0,
hypertexnames=false,
colorlinks=true,% Set to false to disable coloring links
citecolor=magenta,% The color of citations
linkcolor=black,% The color of references to document elements (sections, figures, etc)
urlcolor=mdtRed,% The color of hyperlinks (URLs)
pdfstartview={FitV},
unicode,
breaklinks=true,
}
\usepackage{tabularray}
% Used for displaying a sample figure. If possible, figure files should
% be included in EPS format.
%
% If you use the hyperref package, please uncomment the following line
% to display URLs in blue roman font according to Springer's eBook style:
% \renewcommand\UrlFont{\color{blue}\rmfamily}

\begin{document}
%
\title{Supplementary material for TaylorPDENet: Learning PDEs from non-grid Data}
%
%\titlerunning{Abbreviated paper title}
% If the paper title is too long for the running head, you can set
% an abbreviated paper title here
%


% \author{Paul Heinisch \and
% Andrzej Dulny \and
% Andreas Hotho \and
% Anna Krause}
% %
% \authorrunning{P. Heinisch et al.}
% % First names are abbreviated in the running head.
% % If there are more than two authors, 'et al.' is used.
% %
% \institute{University of Würzburg, Germany \\
% \email{paul.heinisch@stud-mail.uni-wuerzburg.de}\\
% \email{\{dulny,andreas.hotho,anna.krause\}@uni-wuerzburg.de}
% }
% %

\author{Anonymous Author}
%
\authorrunning{Anonymous}
% First names are abbreviated in the running head.
% If there are more than two authors, 'et al.' is used.
%
\institute{***}



\maketitle              % typeset the header of the contribution
%

%
%
%
\appendix
\section{Hyperparameters}
The hyperparameter study was conducted with Optuna \cite{akiba2019optuna} and led to the following results. During the study, the models were trained with $L=20$ prediction steps, on a grid of $64 \times 64$ or on non-grid data $64 ^2$ downsampled points, and $10$ epochs. For the TaylorPDENet $16$ neighbors were used. The study was conducted with the TPESampler and 100 trials.
\begin{center}
\begin{tblr}{m{55pt}@{\hskip 0.5in}l}
    \textbf{ \\ TaylorPDENet \\ (non-grid)} & 
        \begin{tabular}{l@{\hskip 0.4in}l@{\hskip 0.4in}l}
        \textbf{Parameter} & \textbf{Result} & \textbf{Search space} \\
        \\
        Weight Decay & $1e^{-7}$ & log$[1e^{-3}, 1e^{-8}]$\\
        Learning Rate & $1e^{-3}$ & log$[1e^{-2}, 1e^{-4}]$\\
        Batch Size & $9$ & $[1, 16]$\\
        Neighbors & $20$ & $[10, 30]$\\
        \end{tabular} \\
        \hline
    \textbf{TaylorPDENet \\ (grid)} & 
        \begin{tabular}{l@{\hskip 0.4in}l@{\hskip 0.57in}l}
        Weight Decay & $6e^{-5}$ & log$[1e^{-3}, 1e^{-8}]$\\
        Learning Rate & $8e^{-3}$ & log$[1e^{-2}, 1e^{-4}]$\\
        Batch Size & $9$ & $[1, 16]$\\
        Neighbors & $28$ & $[10, 30]$\\
        \end{tabular}\\
        \hline
    \textbf{PDE-Net} & 
        \begin{tabular}{l@{\hskip 0.4in}l@{\hskip 0.35in}l}
        Weight Decay & $1e^{-6}$ & log$[1e^{-3}, 1e^{-8}]$\\
        Learning Rate & $1e^{-3}$ & log$[1e^{-2}, 1e^{-4}]$\\
        Batch Size & $2$ & $[1, 16]$\\
        Constraint & $moment$ & $[moment, frozen]$\\
        Kernel Size & $5$ & $[3, 5, 7]$\\
        \end{tabular}
\end{tblr}\\
\end{center}
In the table above the hyperparameters can be found as well with their search spaces. The $log$ in front of the search space indicates a logarithmic search space.
These hyperparameters were used for the most part during the experiments, except when stated otherwise.

\section{Additional experiments}
In this section, we are showing some supplementary results. First, we show results for the TaylorPDENet on downsampled data with the size $32^2$. Second, we will conduct an experiment investigating the relationship between the number of prediction steps and reconstruction and forecasting MSE.

\begin{figure}[!h]
% focus on non-grid and error map
    \centering
    \begin{subfigure}{0.23\textwidth}
        \centering
        \includegraphics[width=\textwidth]{Figures/e5_256_p32_p.png}
    \end{subfigure}
    \hfill
    \begin{subfigure}{0.75\textwidth}
        \centering
        \includegraphics[width=\textwidth]{Figures/e5_256_p32.png}
    \end{subfigure}
\caption{Prediction after $150$ prediction steps for a PDE learned from irregularly spaced data points with only $32^2$ sampled points.}
\label{fig:downsampled_points}
\end{figure}

\subsection{Downsampling}
In \Cref{fig:downsampled_points,fig:32nongrid} results can be seen for the TalorPDENet on $32^2$ points. The effectiveness, even on sparse data, can be seen.

\begin{figure}[!ht]
    \centering
    \includegraphics[width=\textwidth]{Figures/e5_256_p32_2.png}
    \caption{Prediction with TaylorPDENet on $32^2$ downsampled points after an increasing number of prediction steps.}
    \label{fig:32nongrid}
\end{figure}
%
\begin{figure}[!ht]
    \centering
    \includegraphics[width=\textwidth]{Figures/e2_64_2.png}
    \caption{Prediction with TaylorPDENet on $64 \times 64$ points after an increasing number of prediction steps.}
    \label{fig:64grid}
\end{figure}
%
\subsection{Used prediction steps during training}
With this experiment, we show the influence of different numbers of prediction steps. In \Cref{fig:dt_reconstruction} the influence of the number of prediction steps during training for the reconstruction can be seen. The reconstruction MSE shows that the final result is in the same order of magnitude. Furthermore, the model is tested with $L=150$ prediction steps as well and it can be seen that the result are very similar. Due to the very similar reconstruction MSE, the prediction is very similar as well.

\begin{table}[!h]
\label{tab:dt_prediction}
\centering
\begin{tabular}{l@{\hskip 0.1in}c@{\hskip 0.1in}c@{\hskip 0.1in}c@{\hskip 0.1in}c@{\hskip 0.1in}c@{\hskip 0.1in}c} %@{\hskip 0.1in}
% e3_256
\toprule
Prediction steps & 1 & 2 & 4 & 8 & 16 & 32\\
\midrule
Forecasting MSE ($10^{-4}$) & $1.07$ & $1.07$ & $1.07$ & $1.06$ & $1.04$ & $1.02$ \\
\bottomrule
\end{tabular}
\caption{The forecasting MSE for $L=150$ prediction steps. The model was trained with the indicated number of prediction steps on non-grid data.}
\end{table}

\begin{figure}[!ht]
    \centering
    \includegraphics[width=\textwidth]{Figures/reconstruct_loss.png}
    \caption{The reconstruction MSE during training. It can be seen that the reconstruction MSE is similar at the end of the training.}
    \label{fig:dt_reconstruction}
\end{figure}

%
% ---- Bibliography ----
%
% BibTeX users should specify bibliography style 'splncs04'.
% References will then be sorted and formatted in the correct style.
%
%\clearpage
\bibliographystyle{splncs04}
\bibliography{bibliography}
%
\end{document}